% D2.2 Imitation learning: demonstrations -> dataset -> policy -> robot.
% Shared figure colours (colour only in figures).
\definecolor{robotfill}{HTML}{EAF2FC}\definecolor{robotline}{HTML}{2A78D6}
\definecolor{serverfill}{HTML}{FDEEE7}\definecolor{serverline}{HTML}{EB6834}
\definecolor{netfill}{HTML}{F0EFEC}\definecolor{greenfill}{HTML}{E9F5EC}\definecolor{greenline}{HTML}{2F9E58}
\definecolor{purplefill}{HTML}{F1ECF9}\definecolor{purpleline}{HTML}{7A55C7}
\definecolor{ink}{HTML}{0B0B0B}\definecolor{ink2}{HTML}{52514E}\definecolor{faultc}{HTML}{C9C6BF}

\begin{tikzpicture}[font=\small, node distance=5mm and 9mm,
  st/.style={draw=#1, thick, rounded corners=3pt, fill=white, align=center, text width=31mm, minimum height=15mm},
  arr/.style={-{Latex[length=2.4mm]}, very thick, ink2}]
\node[st=robotline] (demo) {{\bfseries Demonstration}\\\footnotesize a person moves the leader arm; the follower copies it};
\node[st=robotline, right=of demo] (rec) {{\bfseries Recording}\\\footnotesize images, joint states and actions at 30~fps};
\node[st=serverline, right=of rec] (train) {{\bfseries Training}\\\footnotesize the policy learns observation
  $\rightarrow$ action chunk};
\node[st=greenline, below=12mm of train] (pol) {{\bfseries Policy}\\\footnotesize new observation in, a chunk of actions
  out};
\node[st=robotline, left=of pol] (robot) {{\bfseries Robot}\\\footnotesize executes the actions; the cameras see the
  result};
\draw[arr] (demo) -- (rec); \draw[arr] (rec) -- node[above, font=\scriptsize]{dataset} (train);
\draw[arr] (train) -- node[right, font=\scriptsize]{weights} (pol);
\draw[arr] ([yshift=2mm]pol.west) -- node[above, font=\scriptsize]{actions} ([yshift=2mm]robot.east);
\draw[arr] ([yshift=-2mm]robot.east) -- node[below, font=\scriptsize]{observation} ([yshift=-2mm]pol.west);
\end{tikzpicture}
